% Prezentacja artykułu bazowego (Lab 4, 7.11): 8 slajdów, do 5 minut.
% Kompilacja: latexmk -pdf base_article.tex
\documentclass[aspectratio=169,11pt]{beamer}
\usepackage[T1]{fontenc}
\usepackage[utf8]{inputenc}
\usepackage{polski}
\usepackage{lmodern}
\usepackage{booktabs}
\usepackage{hyperref}

\usetheme{default}
\usecolortheme{seahorse}
\setbeamertemplate{navigation symbols}{}
\setbeamertemplate{footline}[frame number]
\setbeamerfont{footnote}{size=\tiny}

\title{Benchmarking OCR and vision-language models\\for Turkish text recognition}
\subtitle{Artykuł bazowy: Y.~Yılmaz, E.\,G.~Hanoğlu, A.\,G.~Özkan, K.~Öztoprak\\
  \emph{International Journal on Document Analysis and Recognition}, 2026,
  \href{https://doi.org/10.1007/s10032-026-00613-6}{doi:10.1007/s10032-026-00613-6}}
\author{[Imię Nazwisko], [Imię Nazwisko], [Imię Nazwisko]}
\date{Laboratorium, 7 listopada 2026}

\begin{document}

\begin{frame}
  \titlepage
\end{frame}

% 1 --------------------------------------------------------------------------
\begin{frame}{Czasopismo i autorzy}
  \begin{columns}[T]
    \begin{column}{0.5\textwidth}
      \textbf{IJDAR}, Springer, ISSN 1433-2833 / 1433-2825
      \begin{itemize}
        \item wykaz MNiSW: \textbf{100 pkt}, m.in. informatyka techniczna i~telekomunikacja
          oraz informatyka\footnotemark[1]
        \item Impact Factor 2025: \textbf{1,7} (bez autocytowań 1,6)\footnotemark[2]
        \item SJR 2025: 0,440, \textbf{Q2}; indeks h czasopisma: 64\footnotemark[3]
        \item typ: \emph{original paper} (nowe badania)
      \end{itemize}
    \end{column}
    \begin{column}{0.48\textwidth}
      \small
      Konya Food and Agriculture University, Turcja\\[4pt]
      \begin{tabular}{lr}
        \toprule
        Autor & indeks h\footnotemark[4] \\
        \midrule
        Yasin Yılmaz & b.d.\footnotemark[5] \\
        Erol Görkem Hanoğlu & 0 \\
        Ayşe Gül Özkan & 1 \\
        Kasım Öztoprak & 11 \\
        \bottomrule
      \end{tabular}
    \end{column}
  \end{columns}
  \footnotetext[1]{Wykaz czasopism naukowych MNiSW, styczeń 2024.}
  \footnotetext[2]{Clarivate, Journal Citation Reports.}
  \footnotetext[3]{SCImago Journal \& Country Rank.}
  \footnotetext[4]{OpenAlex, stan na 22.10.2026.}
  \footnotetext[5]{ORCID 0009-0008-3985-4624 bez profilu w OpenAlex; popularne nazwisko uniemożliwia jednoznaczne przypisanie.}
\end{frame}

% 2 --------------------------------------------------------------------------
\begin{frame}{Cel, słowa kluczowe, nowość}
  \textbf{Słowa kluczowe:} OCR, VLMs, Turkish text recognition, synthetic dataset
  \medskip

  \textbf{Cel:} pierwsza systematyczna ocena klasycznego OCR i modeli wizyjno-językowych (VLM)
  w~rozpoznawaniu tekstu tureckiego, języka o~mniejszych zasobach treningowych
  (znaki ç, ğ, ı, İ, ö, ş, ü; język aglutynacyjny, długie słowa).
  \medskip

  \textbf{Elementy nowości:}
  \begin{itemize}
    \item syntetyczny zbiór 13\,200 obrazów z~kontrolowanymi kategoriami trudności,
    \item 10 modeli trzech typów porównanych w~jednym protokole,
    \item analiza wpływu znaków tureckich, długości słowa, kontekstu zdania i~zniekształceń,
    \item otwarte dane (Zenodo, CC BY 4.0) i~kod (GitHub).
  \end{itemize}
\end{frame}

% 3 --------------------------------------------------------------------------
\begin{frame}{Materiał}
  \begin{itemize}
    \item Teksty: \emph{Nutuk} Atatürka (6600 obrazów) i~turecka Wikipedia (6600) oraz 90 zdjęć rzeczywistych
      (\emph{Gold Standard}).
    \item Trzy typy tekstu: \textbf{drukowany} i~\textbf{odręczny} z~generatora TRDG 1.8.0
      (13 krojów drukowanych, 340 fontów odręcznych z~Google Fonts) oraz \textbf{tekst w~scenie}
      nakładany własnymi skryptami na tła COCO.
    \item Cztery kategorie, 600 obrazów na podkategorię (po 200 na typ tekstu):
      \begin{enumerate}
        \item ze znakami tureckimi / bez nich,
        \item słowa krótkie / długie (próg 10 znaków),
        \item pojedyncze słowo / zdanie,
        \item zniekształcenia: czysty, obrót, szum, rozmycie, niska rozdzielczość.
      \end{enumerate}
  \end{itemize}
\end{frame}

% 4 --------------------------------------------------------------------------
\begin{frame}{Metody}
  \begin{itemize}
    \item \textbf{Klasyczne OCR:} TrOCR (large-printed), PaddleOCR.
    \item \textbf{Wyspecjalizowany model OCR:} GOT-OCR2.
    \item \textbf{Otwarte VLM:} Qwen2.5-VL-7B, InternVL3-8B, DeepSeek-VL2-tiny, Phi-3.5-Vision,
      Pixtral-12B.
    \item \textbf{Komercyjne VLM:} GPT-4o, Claude.
    \item Ten sam prompt dla VLM: \emph{,,Extract the Turkish text from the image exactly as it appears\ldots''};
      74 odpowiedzi (0,56\%) z~dodanym komentarzem powtórzono z~promptem zapasowym.
    \item Miary: \textbf{CER} (główna) i~\textbf{WER} (zdania), biblioteka jiwer; małe litery
      (TrOCR zwraca same wielkie), CER obcięty do 1,0; przedziały ufności 95\% metodą bootstrap.
  \end{itemize}
\end{frame}

% 5 --------------------------------------------------------------------------
\begin{frame}{Wyniki}
  \begin{columns}[T]
    \begin{column}{0.52\textwidth}
      \begin{itemize}
        \item Najlepsze: \textbf{GPT-4o} i~\textbf{Qwen2.5-VL-7B}; otwarty Qwen dorównuje komercyjnemu modelowi.
        \item Najgorsze: GOT-OCR2 (0,391) i~TrOCR (0,383); u~autorów ostatni TrOCR, bo liczą CER bez spacji (GOT-OCR2 gubi spacje).
        \item Znaki tureckie pogarszają wynik wszystkich modeli poza GPT-4o (0,055 / 0,065).
        \item Długość słowa nie ma dużego wpływu; najtrudniejsze są \textbf{pismo odręczne}
          i~\textbf{niska rozdzielczość} (największy wpływ u~9 z~10 modeli).
      \end{itemize}
    \end{column}
    \begin{column}{0.45\textwidth}
      \footnotesize
      Średni CER (przeliczony z~ich danych, 6600 obrazów):\\[3pt]
      \begin{tabular}{lr}
        \toprule
        Model & CER \\
        \midrule
        GPT-4o & 0,070 \\
        Qwen2.5-VL-7B & 0,077 \\
        InternVL3-8B & 0,182 \\
        DeepSeek-VL2-tiny & 0,191 \\
        Phi-3.5-Vision & 0,238 \\
        PaddleOCR & 0,310 \\
        Claude & 0,349 \\
        Pixtral-12B & 0,370 \\
        TrOCR & 0,383 \\
        GOT-OCR2 & 0,391 \\
        \bottomrule
      \end{tabular}
    \end{column}
  \end{columns}
\end{frame}

% 6 --------------------------------------------------------------------------
\begin{frame}{Odtwarzalność: co sprawdziliśmy}
  Dane (Zenodo, CC BY 4.0) i~kod (GitHub, bez licencji) są dostępne, a~repozytorium zawiera
  predykcje każdego modelu dla każdego obrazu.
  \begin{itemize}
    \item \textbf{Przeliczenie:} z~ich predykcji odtwarzamy większość liczb z~tekstu (TrOCR 0,380
      wobec 0,38; GPT-4o na czystych obrazach 0,017 wobec 0,017), ale nie wszystkie
      (PaddleOCR na czystych obrazach 0,260 wobec 0,239).
    \item \textbf{Ponowne uruchomienie} TrOCR i~PaddleOCR na 330 obrazach (10 na komórkę):
      po \textbf{99,7\% predykcji identycznych} z~predykcjami autorów.
    \item \textbf{Rozbieżności:}
      \begin{itemize}
        \item kod i~opublikowany JSON liczą CER jako odległość / długość \emph{predykcji}, a~nie prawdy;
        \item liczby z~tekstu odpowiadają CER \emph{bez spacji i~interpunkcji}, czego artykuł nie podaje;
        \item opublikowane wyniki dotyczą Claude 3.7 Sonnet (CER 0,349), w~tekście 0,197;
        \item tekstury tła papieru nie zostały udostępnione.
      \end{itemize}
  \end{itemize}
\end{frame}

% 7 --------------------------------------------------------------------------
\begin{frame}{Ograniczenia i nasz wkład}
  \textbf{Ograniczenia (wg autorów):} dane syntetyczne (pismo odręczne z~fontów, tekst w~scenie
  na tłach COCO), mały zbiór rzeczywisty, tylko język turecki.
  \medskip

  \textbf{Nasze badanie:} ta sama metodologia dla \textbf{języka polskiego} (ą, ć, ę, ł, ń, ó, ś, ź, ż):
  \begin{itemize}
    \item zbiór syntetyczny z~tekstów Wolnych Lektur w~domenie publicznej (m.in. \emph{Lalka} B.~Prusa):
      druk i~pismo odręczne, te same kategorie;
    \item modele: Qwen2.5-VL, PaddleOCR, TrOCR oraz \textbf{Tesseract (pol)} i~\textbf{EasyOCR (pl)},
      czyli klasyczne OCR z~modelem polskim, których autorzy nie badali;
    \item poprawiona miara CER, otwarty kod i~dane.
  \end{itemize}
\end{frame}

\end{document}
